\chapter{Citações e Referências}

Citações e referências são parte central da integridade acadêmica. Elas mostram de onde vêm informações, conceitos, dados e argumentos utilizados no trabalho. Também permitem que o leitor verifique as fontes e compreenda o diálogo do TCC com a produção existente.

\section{Quando citar}

Deve-se citar sempre que o texto utilizar ideia, conceito, dado, interpretação, trecho literal, imagem, tabela, método, classificação ou argumento de outra fonte. A citação é necessária mesmo quando o estudante reescreve a ideia com suas próprias palavras.

Não é necessário citar conhecimentos amplamente consolidados e gerais, mas, em caso de dúvida, é mais seguro indicar a fonte que sustentou a afirmação.

\section{Citação indireta}

A citação indireta apresenta a ideia de outro autor com as palavras do estudante. Ela deve preservar o sentido original e indicar a fonte. Parafrasear não significa trocar algumas palavras; significa compreender a ideia e reescrevê-la de modo próprio.

\section{Citação direta curta}

A citação direta curta reproduz literalmente trecho de uma fonte e permanece incorporada ao parágrafo, entre aspas, com indicação de autoria, ano e página ou outro localizador disponível. Deve ser usada com moderação, apenas quando a formulação original for importante.

\section{Citação direta longa}

A citação direta longa, com mais de três linhas, é destacada do parágrafo, sem aspas, em espaço simples e com fonte menor. A ABNT recomenda recuo de 4 cm; a UnDF adota esse valor como decisão institucional validada no âmbito do projeto. A citação deve indicar página ou outro localizador disponível e ser comentada pelo estudante. Não se deve inserir bloco citado sem explicar sua relevância para o argumento.

\section{Citação de citação}

A citação de citação deve ser excepcional e evitada quando a fonte original puder ser consultada. Na chamada, identifica-se a autoria original e a fonte efetivamente consultada com a expressão \emph{apud}; na lista de referências, registra-se somente a obra consultada. O ideal é que o estudante consulte diretamente as obras centrais de seu tema.

\section{Localizadores e documentos não paginados}

Nas citações diretas, a página deve ser indicada quando existir. Em documentos sem paginação, use seção, capítulo, artigo, parágrafo, posição, timestamp ou outro identificador capaz de levar o leitor ao trecho. Na citação indireta, a indicação de página ou localização é opcional, mas pode ser útil em obras extensas.

\section{Sistema autor-data}

Nas chamadas entre parênteses, sobrenomes de pessoas e nomes de entidades são grafados em letras maiúsculas e minúsculas, por exemplo, (Silva, 2024) e (Associação Brasileira de Normas Técnicas, 2024). Siglas consolidadas permanecem em caixa alta, como (IBGE, 2025). A abreviação \emph{et al.} pode ser utilizada em obras com quatro ou mais autores e deve seguir padrão uniforme em todo o trabalho.

\section{Referências}

A lista de referências deve permitir identificação e localização das fontes. Obras impressas, artigos, capítulos, legislação, documentos digitais, sites, relatórios, vídeos e outros materiais exigem informações específicas.

Para situações especiais de chamada e localização, consulte o Capítulo \ref{chap:casos-citacao}. Para modelos completos de referências por tipo documental, consulte o Capítulo \ref{chap:modelos-referencias}.

Erros comuns em referências incluem:

\begin{itemize}
  \item autor citado no texto ausente da lista final;
  \item referência listada sem citação no texto;
  \item links quebrados ou sem data de acesso quando necessária;
  \item títulos incompletos;
  \item autores grafados de formas diferentes;
  \item mistura de estilos;
  \item ausência de informações essenciais.
\end{itemize}

\caixaatencao{A conferência cruzada é obrigatória: toda fonte citada deve aparecer nas referências, e toda referência deve corresponder a uma fonte efetivamente utilizada no trabalho.}

\section{Fontes confiáveis}

Nem toda fonte é adequada para TCC. O estudante deve priorizar literatura científica, livros acadêmicos, documentos oficiais, legislação, relatórios técnicos e bases reconhecidas. Fontes jornalísticas, blogs e conteúdos de redes sociais podem ser úteis em objetos específicos, mas precisam de justificativa e análise crítica.
